% Раздел для включения через % Раздел для включения через % Раздел для включения через % Раздел для включения через \input{adp_matrix_statistics.tex}
% Требуемые пакеты: amsmath, amssymb

\section{Матричное вычисление статистик Average Derivative Procedure}
\label{sec:adp-matrix-statistics}

Рассматриваются наблюдения
\[
X_i\in\mathbb R^d,\qquad Y_i\in\mathbb R,
\qquad i=1,\ldots,n,
\]
и набор локальных центров с индексами $j=1,\ldots,J$. Для каждого центра заданы веса
\[
w_{ji}=K\!\left(\left\|T(X_i-x_j)\right\|_2^2\right)\geq 0
\]
и $P=n_{\Phi}$ единичных направлений
$\phi_{jr}\in\mathbb R^d$, $r=1,\ldots,P$.
В ADP центрирование выполняется относительно локального взвешенного среднего
\begin{equation}
M_j=
\frac{\sum_{i=1}^n w_{ji}X_i}
     {\sum_{i=1}^n w_{ji}}.
\label{eq:adp-local-mean}
\end{equation}
Поэтому первый взвешенный момент равен нулю:
\begin{equation}
\sum_{i=1}^n w_{ji}(X_i-M_j)=0.
\label{eq:adp-zero-first-moment}
\end{equation}
Для каждого направления вычисляются статистики
\begin{align}
I_{jr}
&=
\sum_{i=1}^n
Y_i\,\phi_{jr}^{\mathsf T}(X_i-M_j)\,w_{ji},
\label{eq:adp-I-direct}
\\
U_{jr}
&=
\sum_{i=1}^n
(X_i-M_j)\,\phi_{jr}^{\mathsf T}(X_i-M_j)\,w_{ji}
\in\mathbb R^d.
\label{eq:adp-U-direct}
\end{align}
Строки $U_{jr}^{\mathsf T}$ образуют матрицу
$U_j\in\mathbb R^{P\times d}$, а значения $I_{jr}$ образуют
$I_j\in\mathbb R^P$. Локальная линейная модель приводит к приближению
\[
I_j\approx f'_jU_j\beta,
\qquad \|\beta\|_2=1.
\]
Свободный локальный коэффициент отсутствует именно благодаря
\eqref{eq:adp-zero-first-moment}.

\subsection{Плотная матричная факторизация}

Матрицу наблюдений будем хранить строками:
\[
X=
\begin{pmatrix}
X_1^{\mathsf T}\\[-1mm]
\vdots\\[-1mm]
X_n^{\mathsf T}
\end{pmatrix}
\in\mathbb R^{n\times d},
\qquad
y=(Y_1,\ldots,Y_n)^{\mathsf T}.
\]
Для батча из $B$ центров введём матрицу весов
$W\in\mathbb R^{B\times n}$ и вектор локальных масс
\begin{equation}
m=W\mathbf 1_n,
\qquad m_j=\sum_iw_{ji}.
\label{eq:adp-mass}
\end{equation}
Строки с $m_j\leq m_{\min}$ должны исключаться либо пересчитываться с большим bandwidth.
Нормированные веса задаются матрицей
\begin{equation}
A=\operatorname{diag}(m)^{-1}W,
\qquad A\mathbf 1_n=\mathbf 1_B.
\label{eq:adp-normalized-weights}
\end{equation}
Матрица локальных средних получается одним матричным умножением:
\begin{equation}
\mathcal M=AX\in\mathbb R^{B\times d},
\qquad \mathcal M_{j:}=M_j^{\mathsf T}.
\label{eq:adp-means-matmul}
\end{equation}

Направления для батча хранятся в тензоре
$\Phi\in\mathbb R^{B\times P\times d}$.
Сначала вычисляются проекции всех наблюдений:
\begin{equation}
R=\operatorname{bmm}(\Phi,X^{\mathsf T})
\in\mathbb R^{B\times P\times n},
\qquad
R_{jri}=\phi_{jr}^{\mathsf T}X_i.
\label{eq:adp-raw-projections}
\end{equation}
Проекции локальных средних имеют вид
\begin{equation}
c_{jr}=\phi_{jr}^{\mathsf T}M_j.
\label{eq:adp-projected-mean}
\end{equation}
После broadcasting получаем центрированные проекции
\begin{equation}
Q_{jri}=R_{jri}-c_{jr}
       =\phi_{jr}^{\mathsf T}(X_i-M_j).
\label{eq:adp-Q}
\end{equation}
Главный вспомогательный тензор строится произведением Адамара:
\begin{equation}
H=Q\odot A_{:\,,\mathrm{None},:},
\qquad
H_{jri}=a_{ji}Q_{jri}.
\label{eq:adp-H}
\end{equation}
Тогда обе статистики вычисляются через матричные произведения:
\begin{align}
I
&=
 m_{:\,,\mathrm{None}}\odot(Hy)
\in\mathbb R^{B\times P},
\label{eq:adp-I-matrix}
\\
U
&=
 m_{:\,,\mathrm{None},\mathrm{None}}
 \odot(HX)
\in\mathbb R^{B\times P\times d}.
\label{eq:adp-U-matrix}
\end{align}
Формула \eqref{eq:adp-U-matrix} следует из
\[
\sum_iw_{ji}Q_{jri}M_j
=M_j\phi_{jr}^{\mathsf T}\sum_iw_{ji}(X_i-M_j)=0.
\]
Следовательно, тензор разностей размера $B\times n\times d$ формировать не требуется.
Основной временный объект имеет размер $B\times P\times n$; при $P\ll d$
это существенно уменьшает память.

\paragraph{Устойчивая форма.}
Для защиты от ошибок округления рекомендуется предварительно глобально центрировать данные:
\[
X^{c}=X-\mathbf 1_nx_0^{\mathsf T},
\qquad
x_0=n^{-1}\sum_iX_i.
\]
Локальные средние вычисляются как $\mathcal M^{c}=AX^{c}$.
После построения $Q$ полезно повторно занулить его взвешенное среднее:
\begin{equation}
\delta_{jr}=\sum_i a_{ji}Q_{jri},
\qquad
Q_{jri}\leftarrow Q_{jri}-\delta_{jr}.
\label{eq:adp-Q-correction}
\end{equation}
Пусть $s=H\mathbf 1_n$ и $\bar y=Ay$. Численно защищённые формулы имеют вид
\begin{align}
I
&=
 m_{:\,,\mathrm{None}}
 \odot\left(Hy-s\odot\bar y\right),
\label{eq:adp-I-stable}
\\
U
&=
 m_{:\,,\mathrm{None},\mathrm{None}}
 \odot\left(
 HX^{c}-s_{:\,,\mathrm{None}}\odot
 \mathcal M^{c}_{:\,,\mathrm{None},:}
 \right).
\label{eq:adp-U-stable}
\end{align}
В точной арифметике $s=0$, поэтому
\eqref{eq:adp-I-stable}--\eqref{eq:adp-U-stable}
совпадают с \eqref{eq:adp-I-matrix}--\eqref{eq:adp-U-matrix}.
В плавающей арифметике дополнительные члены подавляют ошибку,
возникающую при умножении малого остатка $s$ на большой общий сдвиг $X$ или $y$.

\subsection{Батчевая схема}

Полную матрицу весов $J\times n$ и полный тензор проекций
$J\times P\times n$ хранить необязательно. Центры обрабатываются блоками
$j=j_0,\ldots,j_1-1$ размера $B$. Внутри каждого блока выполняются следующие шаги:
\begin{enumerate}
\item вычислить $W_b$, $m_b$, $A_b$ и $\mathcal M_b$;
\item вычислить $R_b=\operatorname{bmm}(\Phi_b,X^{\mathsf T})$;
\item получить $Q_b$, применить коррекцию \eqref{eq:adp-Q-correction};
\item выполнить $H_b=Q_b\odot A_b$;
\item вычислить $I_b$ и $U_b$ по
      \eqref{eq:adp-I-stable}--\eqref{eq:adp-U-stable};
\item записать результат и перейти к следующему блоку.
\end{enumerate}
Если число направлений велико, внутри блока центров вводится дополнительный блок
$r=r_0,\ldots,r_1-1$. Рабочая память плотной версии составляет
\[
O(BPn+BPd+Bn+Bd),
\]
а вычислительная сложность ---
\[
O(BPnd).
\]
Размер батча выбирается по памяти тензора $Q$ и результата $U$.

\subsection{Разреженная матрица весов}

Для ядра с компактным носителем большинство весов равно нулю.
Обозначим
\[
\mathcal N_j=\{i:w_{ji}\neq0\},
\qquad
E=\operatorname{nnz}(W),
\qquad
\bar k=E/J.
\]
На CPU матрицу $W$ целесообразно хранить в формате CSR.
Вектор масс и локальные средние вычисляются sparse--dense операциями:
\begin{equation}
m=W\mathbf 1_n,
\qquad
\mathcal M=\operatorname{diag}(m)^{-1}WX.
\label{eq:adp-sparse-mean}
\end{equation}
Для каждого ненулевого ребра $(j,i)$ и направления $r$ вычисляется только
\begin{equation}
q_{jri}=\phi_{jr}^{\mathsf T}(X_i-M_j),
\qquad i\in\mathcal N_j.
\label{eq:adp-sparse-q}
\end{equation}
Определим разреженную матрицу $H^{(r)}\in\mathbb R^{B\times n}$
с тем же паттерном ненулевых элементов, что у $W$:
\begin{equation}
H^{(r)}_{ji}=
\begin{cases}
\displaystyle \frac{w_{ji}}{m_j}q_{jri}, & i\in\mathcal N_j,\\[2mm]
0, & i\notin\mathcal N_j.
\end{cases}
\label{eq:adp-sparse-H}
\end{equation}
Тогда
\begin{align}
I_{:r}
&=
m\odot\left(H^{(r)}y-s^{(r)}\odot\bar y\right),
\label{eq:adp-sparse-I}
\\
U_{:r:}
&=
m_{:\,,\mathrm{None}}\odot
\left(
H^{(r)}X^{c}
-s^{(r)}_{:\,,\mathrm{None}}\odot\mathcal M^{c}
\right),
\label{eq:adp-sparse-U}
\\
s^{(r)}&=H^{(r)}\mathbf 1_n.
\end{align}
Операции $H^{(r)}y$ и $H^{(r)}X^{c}$ представляют соответственно
sparse matrix--vector и sparse matrix--dense matrix multiplication.
Не следует заранее хранить все $P$ матриц $H^{(r)}$: направления лучше
обрабатывать блоками, повторно используя один CSR-паттерн и меняя только массив значений.

Сложность разреженного варианта равна
\[
O(EPd)=O(J\bar kPd),
\]
а память для статистик и промежуточных значений ---
\[
O(EP+JPd+Jd).
\]
По сравнению с плотной схемой ожидаемое сокращение арифметики имеет порядок
$n/\bar k$.

\subsection{Формат списка соседей для GPU}

Если для каждого центра хранится не более $K$ соседей, разреженную структуру
удобно представить массивами
\[
N\in\mathbb N^{B\times K},
\qquad
W_{\mathrm{loc}}\in\mathbb R^{B\times K}.
\]
Здесь $N_{jk}$ --- индекс $k$-го соседа, а отсутствующие элементы маскируются
нулевым весом. После gather-операции
\[
X_{\mathcal N}=X[N]\in\mathbb R^{B\times K\times d},
\qquad
y_{\mathcal N}=y[N]\in\mathbb R^{B\times K}
\]
вычисления снова сводятся к batched matmul:
\begin{align}
M_j&=\sum_{k=1}^K a_{jk}X_{N_{jk}},
\\
Q_{jrk}&=\phi_{jr}^{\mathsf T}(X_{N_{jk}}-M_j),
\\
H_{jrk}&=a_{jk}Q_{jrk},
\\
I_{jr}&=m_j\sum_{k=1}^KH_{jrk}(Y_{N_{jk}}-\bar Y_j),
\\
U_{jr}&=m_j\sum_{k=1}^KH_{jrk}(X_{N_{jk}}-M_j).
\end{align}
Последние две редукции реализуются как
\[
(B,P,K)(B,K,1)\longrightarrow(B,P,1),
\qquad
(B,P,K)(B,K,d)\longrightarrow(B,P,d).
\]
Этот формат обычно предпочтительнее CSR на GPU, поскольку имеет регулярные размеры
и использует оптимизированные batched GEMM-ядра. Его недостаток --- padding при сильно
различающемся числе соседей.

\subsection{Численная устойчивость и контроль корректности}

Разреженность не меняет математическую задачу, если нулевые веса возникают из
компактного носителя ядра. Отсечение малых положительных весов или top-$K$ выбор
создают другую статистику. Для контроля следует вычислять долю отброшенной массы
\[
r_j=
\frac{\sum_{i:\,w_{ji}\ \mathrm{discarded}}w_{ji}}
     {\sum_iw_{ji}}
\]
и эффективное число соседей
\begin{equation}
n_{\mathrm{eff},j}
=
\frac{\left(\sum_iw_{ji}\right)^2}{\sum_iw_{ji}^2}
=
\frac{1}{\sum_i a_{ji}^2}.
\label{eq:adp-neff}
\end{equation}
Большое число ненулевых элементов не гарантирует устойчивость, если
$n_{\mathrm{eff},j}$ близко к единице.

Рекомендуемые проверки реализации:
\begin{enumerate}
\item сравнить прямые суммы и матричный вариант в \texttt{float64};
\item проверить инвариантность при $X_i\leftarrow X_i+c$ и
      $Y_i\leftarrow Y_i+c_y$;
\item сравнить разные размеры батча и перестановки соседей;
\item отдельно сравнить одну и ту же матрицу $W$ в плотном и CSR-форматах;
\item контролировать
\[
\eta_{jr}=
\frac{|\sum_iw_{ji}Q_{jri}|}
     {\sum_i|w_{ji}Q_{jri}|+\varepsilon};
\]
\item считать строки с малой массой, малым $n_{\mathrm{eff},j}$ или большим
      $\eta_{jr}$ численно ненадёжными.
\end{enumerate}
Для эталонных тестов следует использовать \texttt{float64}. На GPU допустимо хранение
в \texttt{float32}, но локальные массы, средние и редукции желательно аккумулировать
в повышенной точности. Формулы \eqref{eq:adp-I-stable},
\eqref{eq:adp-U-stable}, \eqref{eq:adp-sparse-I} и
\eqref{eq:adp-sparse-U} предпочтительнее некорректированного умножения $HX$.

\paragraph{Рекомендуемая реализация.}
Прямая версия формул \eqref{eq:adp-I-direct}--\eqref{eq:adp-U-direct}
используется как эталон для малых задач. Основное вычислительное ядро использует
факторизацию через $H$, произведение Адамара и матричные умножения. На CPU веса
хранятся в CSR, а на GPU --- в формате списка соседей. Центры и направления
обрабатываются батчами; данные предварительно центрируются и масштабируются, а
нулевой взвешенный момент проекций контролируется явно.

% Требуемые пакеты: amsmath, amssymb

\section{Матричное вычисление статистик Average Derivative Procedure}
\label{sec:adp-matrix-statistics}

Рассматриваются наблюдения
\[
X_i\in\mathbb R^d,\qquad Y_i\in\mathbb R,
\qquad i=1,\ldots,n,
\]
и набор локальных центров с индексами $j=1,\ldots,J$. Для каждого центра заданы веса
\[
w_{ji}=K\!\left(\left\|T(X_i-x_j)\right\|_2^2\right)\geq 0
\]
и $P=n_{\Phi}$ единичных направлений
$\phi_{jr}\in\mathbb R^d$, $r=1,\ldots,P$.
В ADP центрирование выполняется относительно локального взвешенного среднего
\begin{equation}
M_j=
\frac{\sum_{i=1}^n w_{ji}X_i}
     {\sum_{i=1}^n w_{ji}}.
\label{eq:adp-local-mean}
\end{equation}
Поэтому первый взвешенный момент равен нулю:
\begin{equation}
\sum_{i=1}^n w_{ji}(X_i-M_j)=0.
\label{eq:adp-zero-first-moment}
\end{equation}
Для каждого направления вычисляются статистики
\begin{align}
I_{jr}
&=
\sum_{i=1}^n
Y_i\,\phi_{jr}^{\mathsf T}(X_i-M_j)\,w_{ji},
\label{eq:adp-I-direct}
\\
U_{jr}
&=
\sum_{i=1}^n
(X_i-M_j)\,\phi_{jr}^{\mathsf T}(X_i-M_j)\,w_{ji}
\in\mathbb R^d.
\label{eq:adp-U-direct}
\end{align}
Строки $U_{jr}^{\mathsf T}$ образуют матрицу
$U_j\in\mathbb R^{P\times d}$, а значения $I_{jr}$ образуют
$I_j\in\mathbb R^P$. Локальная линейная модель приводит к приближению
\[
I_j\approx f'_jU_j\beta,
\qquad \|\beta\|_2=1.
\]
Свободный локальный коэффициент отсутствует именно благодаря
\eqref{eq:adp-zero-first-moment}.

\subsection{Плотная матричная факторизация}

Матрицу наблюдений будем хранить строками:
\[
X=
\begin{pmatrix}
X_1^{\mathsf T}\\[-1mm]
\vdots\\[-1mm]
X_n^{\mathsf T}
\end{pmatrix}
\in\mathbb R^{n\times d},
\qquad
y=(Y_1,\ldots,Y_n)^{\mathsf T}.
\]
Для батча из $B$ центров введём матрицу весов
$W\in\mathbb R^{B\times n}$ и вектор локальных масс
\begin{equation}
m=W\mathbf 1_n,
\qquad m_j=\sum_iw_{ji}.
\label{eq:adp-mass}
\end{equation}
Строки с $m_j\leq m_{\min}$ должны исключаться либо пересчитываться с большим bandwidth.
Нормированные веса задаются матрицей
\begin{equation}
A=\operatorname{diag}(m)^{-1}W,
\qquad A\mathbf 1_n=\mathbf 1_B.
\label{eq:adp-normalized-weights}
\end{equation}
Матрица локальных средних получается одним матричным умножением:
\begin{equation}
\mathcal M=AX\in\mathbb R^{B\times d},
\qquad \mathcal M_{j:}=M_j^{\mathsf T}.
\label{eq:adp-means-matmul}
\end{equation}

Направления для батча хранятся в тензоре
$\Phi\in\mathbb R^{B\times P\times d}$.
Сначала вычисляются проекции всех наблюдений:
\begin{equation}
R=\operatorname{bmm}(\Phi,X^{\mathsf T})
\in\mathbb R^{B\times P\times n},
\qquad
R_{jri}=\phi_{jr}^{\mathsf T}X_i.
\label{eq:adp-raw-projections}
\end{equation}
Проекции локальных средних имеют вид
\begin{equation}
c_{jr}=\phi_{jr}^{\mathsf T}M_j.
\label{eq:adp-projected-mean}
\end{equation}
После broadcasting получаем центрированные проекции
\begin{equation}
Q_{jri}=R_{jri}-c_{jr}
       =\phi_{jr}^{\mathsf T}(X_i-M_j).
\label{eq:adp-Q}
\end{equation}
Главный вспомогательный тензор строится произведением Адамара:
\begin{equation}
H=Q\odot A_{:\,,\mathrm{None},:},
\qquad
H_{jri}=a_{ji}Q_{jri}.
\label{eq:adp-H}
\end{equation}
Тогда обе статистики вычисляются через матричные произведения:
\begin{align}
I
&=
 m_{:\,,\mathrm{None}}\odot(Hy)
\in\mathbb R^{B\times P},
\label{eq:adp-I-matrix}
\\
U
&=
 m_{:\,,\mathrm{None},\mathrm{None}}
 \odot(HX)
\in\mathbb R^{B\times P\times d}.
\label{eq:adp-U-matrix}
\end{align}
Формула \eqref{eq:adp-U-matrix} следует из
\[
\sum_iw_{ji}Q_{jri}M_j
=M_j\phi_{jr}^{\mathsf T}\sum_iw_{ji}(X_i-M_j)=0.
\]
Следовательно, тензор разностей размера $B\times n\times d$ формировать не требуется.
Основной временный объект имеет размер $B\times P\times n$; при $P\ll d$
это существенно уменьшает память.

\paragraph{Устойчивая форма.}
Для защиты от ошибок округления рекомендуется предварительно глобально центрировать данные:
\[
X^{c}=X-\mathbf 1_nx_0^{\mathsf T},
\qquad
x_0=n^{-1}\sum_iX_i.
\]
Локальные средние вычисляются как $\mathcal M^{c}=AX^{c}$.
После построения $Q$ полезно повторно занулить его взвешенное среднее:
\begin{equation}
\delta_{jr}=\sum_i a_{ji}Q_{jri},
\qquad
Q_{jri}\leftarrow Q_{jri}-\delta_{jr}.
\label{eq:adp-Q-correction}
\end{equation}
Пусть $s=H\mathbf 1_n$ и $\bar y=Ay$. Численно защищённые формулы имеют вид
\begin{align}
I
&=
 m_{:\,,\mathrm{None}}
 \odot\left(Hy-s\odot\bar y\right),
\label{eq:adp-I-stable}
\\
U
&=
 m_{:\,,\mathrm{None},\mathrm{None}}
 \odot\left(
 HX^{c}-s_{:\,,\mathrm{None}}\odot
 \mathcal M^{c}_{:\,,\mathrm{None},:}
 \right).
\label{eq:adp-U-stable}
\end{align}
В точной арифметике $s=0$, поэтому
\eqref{eq:adp-I-stable}--\eqref{eq:adp-U-stable}
совпадают с \eqref{eq:adp-I-matrix}--\eqref{eq:adp-U-matrix}.
В плавающей арифметике дополнительные члены подавляют ошибку,
возникающую при умножении малого остатка $s$ на большой общий сдвиг $X$ или $y$.

\subsection{Батчевая схема}

Полную матрицу весов $J\times n$ и полный тензор проекций
$J\times P\times n$ хранить необязательно. Центры обрабатываются блоками
$j=j_0,\ldots,j_1-1$ размера $B$. Внутри каждого блока выполняются следующие шаги:
\begin{enumerate}
\item вычислить $W_b$, $m_b$, $A_b$ и $\mathcal M_b$;
\item вычислить $R_b=\operatorname{bmm}(\Phi_b,X^{\mathsf T})$;
\item получить $Q_b$, применить коррекцию \eqref{eq:adp-Q-correction};
\item выполнить $H_b=Q_b\odot A_b$;
\item вычислить $I_b$ и $U_b$ по
      \eqref{eq:adp-I-stable}--\eqref{eq:adp-U-stable};
\item записать результат и перейти к следующему блоку.
\end{enumerate}
Если число направлений велико, внутри блока центров вводится дополнительный блок
$r=r_0,\ldots,r_1-1$. Рабочая память плотной версии составляет
\[
O(BPn+BPd+Bn+Bd),
\]
а вычислительная сложность ---
\[
O(BPnd).
\]
Размер батча выбирается по памяти тензора $Q$ и результата $U$.

\subsection{Разреженная матрица весов}

Для ядра с компактным носителем большинство весов равно нулю.
Обозначим
\[
\mathcal N_j=\{i:w_{ji}\neq0\},
\qquad
E=\operatorname{nnz}(W),
\qquad
\bar k=E/J.
\]
На CPU матрицу $W$ целесообразно хранить в формате CSR.
Вектор масс и локальные средние вычисляются sparse--dense операциями:
\begin{equation}
m=W\mathbf 1_n,
\qquad
\mathcal M=\operatorname{diag}(m)^{-1}WX.
\label{eq:adp-sparse-mean}
\end{equation}
Для каждого ненулевого ребра $(j,i)$ и направления $r$ вычисляется только
\begin{equation}
q_{jri}=\phi_{jr}^{\mathsf T}(X_i-M_j),
\qquad i\in\mathcal N_j.
\label{eq:adp-sparse-q}
\end{equation}
Определим разреженную матрицу $H^{(r)}\in\mathbb R^{B\times n}$
с тем же паттерном ненулевых элементов, что у $W$:
\begin{equation}
H^{(r)}_{ji}=
\begin{cases}
\displaystyle \frac{w_{ji}}{m_j}q_{jri}, & i\in\mathcal N_j,\\[2mm]
0, & i\notin\mathcal N_j.
\end{cases}
\label{eq:adp-sparse-H}
\end{equation}
Тогда
\begin{align}
I_{:r}
&=
m\odot\left(H^{(r)}y-s^{(r)}\odot\bar y\right),
\label{eq:adp-sparse-I}
\\
U_{:r:}
&=
m_{:\,,\mathrm{None}}\odot
\left(
H^{(r)}X^{c}
-s^{(r)}_{:\,,\mathrm{None}}\odot\mathcal M^{c}
\right),
\label{eq:adp-sparse-U}
\\
s^{(r)}&=H^{(r)}\mathbf 1_n.
\end{align}
Операции $H^{(r)}y$ и $H^{(r)}X^{c}$ представляют соответственно
sparse matrix--vector и sparse matrix--dense matrix multiplication.
Не следует заранее хранить все $P$ матриц $H^{(r)}$: направления лучше
обрабатывать блоками, повторно используя один CSR-паттерн и меняя только массив значений.

Сложность разреженного варианта равна
\[
O(EPd)=O(J\bar kPd),
\]
а память для статистик и промежуточных значений ---
\[
O(EP+JPd+Jd).
\]
По сравнению с плотной схемой ожидаемое сокращение арифметики имеет порядок
$n/\bar k$.

\subsection{Формат списка соседей для GPU}

Если для каждого центра хранится не более $K$ соседей, разреженную структуру
удобно представить массивами
\[
N\in\mathbb N^{B\times K},
\qquad
W_{\mathrm{loc}}\in\mathbb R^{B\times K}.
\]
Здесь $N_{jk}$ --- индекс $k$-го соседа, а отсутствующие элементы маскируются
нулевым весом. После gather-операции
\[
X_{\mathcal N}=X[N]\in\mathbb R^{B\times K\times d},
\qquad
y_{\mathcal N}=y[N]\in\mathbb R^{B\times K}
\]
вычисления снова сводятся к batched matmul:
\begin{align}
M_j&=\sum_{k=1}^K a_{jk}X_{N_{jk}},
\\
Q_{jrk}&=\phi_{jr}^{\mathsf T}(X_{N_{jk}}-M_j),
\\
H_{jrk}&=a_{jk}Q_{jrk},
\\
I_{jr}&=m_j\sum_{k=1}^KH_{jrk}(Y_{N_{jk}}-\bar Y_j),
\\
U_{jr}&=m_j\sum_{k=1}^KH_{jrk}(X_{N_{jk}}-M_j).
\end{align}
Последние две редукции реализуются как
\[
(B,P,K)(B,K,1)\longrightarrow(B,P,1),
\qquad
(B,P,K)(B,K,d)\longrightarrow(B,P,d).
\]
Этот формат обычно предпочтительнее CSR на GPU, поскольку имеет регулярные размеры
и использует оптимизированные batched GEMM-ядра. Его недостаток --- padding при сильно
различающемся числе соседей.

\subsection{Численная устойчивость и контроль корректности}

Разреженность не меняет математическую задачу, если нулевые веса возникают из
компактного носителя ядра. Отсечение малых положительных весов или top-$K$ выбор
создают другую статистику. Для контроля следует вычислять долю отброшенной массы
\[
r_j=
\frac{\sum_{i:\,w_{ji}\ \mathrm{discarded}}w_{ji}}
     {\sum_iw_{ji}}
\]
и эффективное число соседей
\begin{equation}
n_{\mathrm{eff},j}
=
\frac{\left(\sum_iw_{ji}\right)^2}{\sum_iw_{ji}^2}
=
\frac{1}{\sum_i a_{ji}^2}.
\label{eq:adp-neff}
\end{equation}
Большое число ненулевых элементов не гарантирует устойчивость, если
$n_{\mathrm{eff},j}$ близко к единице.

Рекомендуемые проверки реализации:
\begin{enumerate}
\item сравнить прямые суммы и матричный вариант в \texttt{float64};
\item проверить инвариантность при $X_i\leftarrow X_i+c$ и
      $Y_i\leftarrow Y_i+c_y$;
\item сравнить разные размеры батча и перестановки соседей;
\item отдельно сравнить одну и ту же матрицу $W$ в плотном и CSR-форматах;
\item контролировать
\[
\eta_{jr}=
\frac{|\sum_iw_{ji}Q_{jri}|}
     {\sum_i|w_{ji}Q_{jri}|+\varepsilon};
\]
\item считать строки с малой массой, малым $n_{\mathrm{eff},j}$ или большим
      $\eta_{jr}$ численно ненадёжными.
\end{enumerate}
Для эталонных тестов следует использовать \texttt{float64}. На GPU допустимо хранение
в \texttt{float32}, но локальные массы, средние и редукции желательно аккумулировать
в повышенной точности. Формулы \eqref{eq:adp-I-stable},
\eqref{eq:adp-U-stable}, \eqref{eq:adp-sparse-I} и
\eqref{eq:adp-sparse-U} предпочтительнее некорректированного умножения $HX$.

\paragraph{Рекомендуемая реализация.}
Прямая версия формул \eqref{eq:adp-I-direct}--\eqref{eq:adp-U-direct}
используется как эталон для малых задач. Основное вычислительное ядро использует
факторизацию через $H$, произведение Адамара и матричные умножения. На CPU веса
хранятся в CSR, а на GPU --- в формате списка соседей. Центры и направления
обрабатываются батчами; данные предварительно центрируются и масштабируются, а
нулевой взвешенный момент проекций контролируется явно.

% Требуемые пакеты: amsmath, amssymb

\section{Матричное вычисление статистик Average Derivative Procedure}
\label{sec:adp-matrix-statistics}

Рассматриваются наблюдения
\[
X_i\in\mathbb R^d,\qquad Y_i\in\mathbb R,
\qquad i=1,\ldots,n,
\]
и набор локальных центров с индексами $j=1,\ldots,J$. Для каждого центра заданы веса
\[
w_{ji}=K\!\left(\left\|T(X_i-x_j)\right\|_2^2\right)\geq 0
\]
и $P=n_{\Phi}$ единичных направлений
$\phi_{jr}\in\mathbb R^d$, $r=1,\ldots,P$.
В ADP центрирование выполняется относительно локального взвешенного среднего
\begin{equation}
M_j=
\frac{\sum_{i=1}^n w_{ji}X_i}
     {\sum_{i=1}^n w_{ji}}.
\label{eq:adp-local-mean}
\end{equation}
Поэтому первый взвешенный момент равен нулю:
\begin{equation}
\sum_{i=1}^n w_{ji}(X_i-M_j)=0.
\label{eq:adp-zero-first-moment}
\end{equation}
Для каждого направления вычисляются статистики
\begin{align}
I_{jr}
&=
\sum_{i=1}^n
Y_i\,\phi_{jr}^{\mathsf T}(X_i-M_j)\,w_{ji},
\label{eq:adp-I-direct}
\\
U_{jr}
&=
\sum_{i=1}^n
(X_i-M_j)\,\phi_{jr}^{\mathsf T}(X_i-M_j)\,w_{ji}
\in\mathbb R^d.
\label{eq:adp-U-direct}
\end{align}
Строки $U_{jr}^{\mathsf T}$ образуют матрицу
$U_j\in\mathbb R^{P\times d}$, а значения $I_{jr}$ образуют
$I_j\in\mathbb R^P$. Локальная линейная модель приводит к приближению
\[
I_j\approx f'_jU_j\beta,
\qquad \|\beta\|_2=1.
\]
Свободный локальный коэффициент отсутствует именно благодаря
\eqref{eq:adp-zero-first-moment}.

\subsection{Плотная матричная факторизация}

Матрицу наблюдений будем хранить строками:
\[
X=
\begin{pmatrix}
X_1^{\mathsf T}\\[-1mm]
\vdots\\[-1mm]
X_n^{\mathsf T}
\end{pmatrix}
\in\mathbb R^{n\times d},
\qquad
y=(Y_1,\ldots,Y_n)^{\mathsf T}.
\]
Для батча из $B$ центров введём матрицу весов
$W\in\mathbb R^{B\times n}$ и вектор локальных масс
\begin{equation}
m=W\mathbf 1_n,
\qquad m_j=\sum_iw_{ji}.
\label{eq:adp-mass}
\end{equation}
Строки с $m_j\leq m_{\min}$ должны исключаться либо пересчитываться с большим bandwidth.
Нормированные веса задаются матрицей
\begin{equation}
A=\operatorname{diag}(m)^{-1}W,
\qquad A\mathbf 1_n=\mathbf 1_B.
\label{eq:adp-normalized-weights}
\end{equation}
Матрица локальных средних получается одним матричным умножением:
\begin{equation}
\mathcal M=AX\in\mathbb R^{B\times d},
\qquad \mathcal M_{j:}=M_j^{\mathsf T}.
\label{eq:adp-means-matmul}
\end{equation}

Направления для батча хранятся в тензоре
$\Phi\in\mathbb R^{B\times P\times d}$.
Сначала вычисляются проекции всех наблюдений:
\begin{equation}
R=\operatorname{bmm}(\Phi,X^{\mathsf T})
\in\mathbb R^{B\times P\times n},
\qquad
R_{jri}=\phi_{jr}^{\mathsf T}X_i.
\label{eq:adp-raw-projections}
\end{equation}
Проекции локальных средних имеют вид
\begin{equation}
c_{jr}=\phi_{jr}^{\mathsf T}M_j.
\label{eq:adp-projected-mean}
\end{equation}
После broadcasting получаем центрированные проекции
\begin{equation}
Q_{jri}=R_{jri}-c_{jr}
       =\phi_{jr}^{\mathsf T}(X_i-M_j).
\label{eq:adp-Q}
\end{equation}
Главный вспомогательный тензор строится произведением Адамара:
\begin{equation}
H=Q\odot A_{:\,,\mathrm{None},:},
\qquad
H_{jri}=a_{ji}Q_{jri}.
\label{eq:adp-H}
\end{equation}
Тогда обе статистики вычисляются через матричные произведения:
\begin{align}
I
&=
 m_{:\,,\mathrm{None}}\odot(Hy)
\in\mathbb R^{B\times P},
\label{eq:adp-I-matrix}
\\
U
&=
 m_{:\,,\mathrm{None},\mathrm{None}}
 \odot(HX)
\in\mathbb R^{B\times P\times d}.
\label{eq:adp-U-matrix}
\end{align}
Формула \eqref{eq:adp-U-matrix} следует из
\[
\sum_iw_{ji}Q_{jri}M_j
=M_j\phi_{jr}^{\mathsf T}\sum_iw_{ji}(X_i-M_j)=0.
\]
Следовательно, тензор разностей размера $B\times n\times d$ формировать не требуется.
Основной временный объект имеет размер $B\times P\times n$; при $P\ll d$
это существенно уменьшает память.

\paragraph{Устойчивая форма.}
Для защиты от ошибок округления рекомендуется предварительно глобально центрировать данные:
\[
X^{c}=X-\mathbf 1_nx_0^{\mathsf T},
\qquad
x_0=n^{-1}\sum_iX_i.
\]
Локальные средние вычисляются как $\mathcal M^{c}=AX^{c}$.
После построения $Q$ полезно повторно занулить его взвешенное среднее:
\begin{equation}
\delta_{jr}=\sum_i a_{ji}Q_{jri},
\qquad
Q_{jri}\leftarrow Q_{jri}-\delta_{jr}.
\label{eq:adp-Q-correction}
\end{equation}
Пусть $s=H\mathbf 1_n$ и $\bar y=Ay$. Численно защищённые формулы имеют вид
\begin{align}
I
&=
 m_{:\,,\mathrm{None}}
 \odot\left(Hy-s\odot\bar y\right),
\label{eq:adp-I-stable}
\\
U
&=
 m_{:\,,\mathrm{None},\mathrm{None}}
 \odot\left(
 HX^{c}-s_{:\,,\mathrm{None}}\odot
 \mathcal M^{c}_{:\,,\mathrm{None},:}
 \right).
\label{eq:adp-U-stable}
\end{align}
В точной арифметике $s=0$, поэтому
\eqref{eq:adp-I-stable}--\eqref{eq:adp-U-stable}
совпадают с \eqref{eq:adp-I-matrix}--\eqref{eq:adp-U-matrix}.
В плавающей арифметике дополнительные члены подавляют ошибку,
возникающую при умножении малого остатка $s$ на большой общий сдвиг $X$ или $y$.

\subsection{Батчевая схема}

Полную матрицу весов $J\times n$ и полный тензор проекций
$J\times P\times n$ хранить необязательно. Центры обрабатываются блоками
$j=j_0,\ldots,j_1-1$ размера $B$. Внутри каждого блока выполняются следующие шаги:
\begin{enumerate}
\item вычислить $W_b$, $m_b$, $A_b$ и $\mathcal M_b$;
\item вычислить $R_b=\operatorname{bmm}(\Phi_b,X^{\mathsf T})$;
\item получить $Q_b$, применить коррекцию \eqref{eq:adp-Q-correction};
\item выполнить $H_b=Q_b\odot A_b$;
\item вычислить $I_b$ и $U_b$ по
      \eqref{eq:adp-I-stable}--\eqref{eq:adp-U-stable};
\item записать результат и перейти к следующему блоку.
\end{enumerate}
Если число направлений велико, внутри блока центров вводится дополнительный блок
$r=r_0,\ldots,r_1-1$. Рабочая память плотной версии составляет
\[
O(BPn+BPd+Bn+Bd),
\]
а вычислительная сложность ---
\[
O(BPnd).
\]
Размер батча выбирается по памяти тензора $Q$ и результата $U$.

\subsection{Разреженная матрица весов}

Для ядра с компактным носителем большинство весов равно нулю.
Обозначим
\[
\mathcal N_j=\{i:w_{ji}\neq0\},
\qquad
E=\operatorname{nnz}(W),
\qquad
\bar k=E/J.
\]
На CPU матрицу $W$ целесообразно хранить в формате CSR.
Вектор масс и локальные средние вычисляются sparse--dense операциями:
\begin{equation}
m=W\mathbf 1_n,
\qquad
\mathcal M=\operatorname{diag}(m)^{-1}WX.
\label{eq:adp-sparse-mean}
\end{equation}
Для каждого ненулевого ребра $(j,i)$ и направления $r$ вычисляется только
\begin{equation}
q_{jri}=\phi_{jr}^{\mathsf T}(X_i-M_j),
\qquad i\in\mathcal N_j.
\label{eq:adp-sparse-q}
\end{equation}
Определим разреженную матрицу $H^{(r)}\in\mathbb R^{B\times n}$
с тем же паттерном ненулевых элементов, что у $W$:
\begin{equation}
H^{(r)}_{ji}=
\begin{cases}
\displaystyle \frac{w_{ji}}{m_j}q_{jri}, & i\in\mathcal N_j,\\[2mm]
0, & i\notin\mathcal N_j.
\end{cases}
\label{eq:adp-sparse-H}
\end{equation}
Тогда
\begin{align}
I_{:r}
&=
m\odot\left(H^{(r)}y-s^{(r)}\odot\bar y\right),
\label{eq:adp-sparse-I}
\\
U_{:r:}
&=
m_{:\,,\mathrm{None}}\odot
\left(
H^{(r)}X^{c}
-s^{(r)}_{:\,,\mathrm{None}}\odot\mathcal M^{c}
\right),
\label{eq:adp-sparse-U}
\\
s^{(r)}&=H^{(r)}\mathbf 1_n.
\end{align}
Операции $H^{(r)}y$ и $H^{(r)}X^{c}$ представляют соответственно
sparse matrix--vector и sparse matrix--dense matrix multiplication.
Не следует заранее хранить все $P$ матриц $H^{(r)}$: направления лучше
обрабатывать блоками, повторно используя один CSR-паттерн и меняя только массив значений.

Сложность разреженного варианта равна
\[
O(EPd)=O(J\bar kPd),
\]
а память для статистик и промежуточных значений ---
\[
O(EP+JPd+Jd).
\]
По сравнению с плотной схемой ожидаемое сокращение арифметики имеет порядок
$n/\bar k$.

\subsection{Формат списка соседей для GPU}

Если для каждого центра хранится не более $K$ соседей, разреженную структуру
удобно представить массивами
\[
N\in\mathbb N^{B\times K},
\qquad
W_{\mathrm{loc}}\in\mathbb R^{B\times K}.
\]
Здесь $N_{jk}$ --- индекс $k$-го соседа, а отсутствующие элементы маскируются
нулевым весом. После gather-операции
\[
X_{\mathcal N}=X[N]\in\mathbb R^{B\times K\times d},
\qquad
y_{\mathcal N}=y[N]\in\mathbb R^{B\times K}
\]
вычисления снова сводятся к batched matmul:
\begin{align}
M_j&=\sum_{k=1}^K a_{jk}X_{N_{jk}},
\\
Q_{jrk}&=\phi_{jr}^{\mathsf T}(X_{N_{jk}}-M_j),
\\
H_{jrk}&=a_{jk}Q_{jrk},
\\
I_{jr}&=m_j\sum_{k=1}^KH_{jrk}(Y_{N_{jk}}-\bar Y_j),
\\
U_{jr}&=m_j\sum_{k=1}^KH_{jrk}(X_{N_{jk}}-M_j).
\end{align}
Последние две редукции реализуются как
\[
(B,P,K)(B,K,1)\longrightarrow(B,P,1),
\qquad
(B,P,K)(B,K,d)\longrightarrow(B,P,d).
\]
Этот формат обычно предпочтительнее CSR на GPU, поскольку имеет регулярные размеры
и использует оптимизированные batched GEMM-ядра. Его недостаток --- padding при сильно
различающемся числе соседей.

\subsection{Численная устойчивость и контроль корректности}

Разреженность не меняет математическую задачу, если нулевые веса возникают из
компактного носителя ядра. Отсечение малых положительных весов или top-$K$ выбор
создают другую статистику. Для контроля следует вычислять долю отброшенной массы
\[
r_j=
\frac{\sum_{i:\,w_{ji}\ \mathrm{discarded}}w_{ji}}
     {\sum_iw_{ji}}
\]
и эффективное число соседей
\begin{equation}
n_{\mathrm{eff},j}
=
\frac{\left(\sum_iw_{ji}\right)^2}{\sum_iw_{ji}^2}
=
\frac{1}{\sum_i a_{ji}^2}.
\label{eq:adp-neff}
\end{equation}
Большое число ненулевых элементов не гарантирует устойчивость, если
$n_{\mathrm{eff},j}$ близко к единице.

Рекомендуемые проверки реализации:
\begin{enumerate}
\item сравнить прямые суммы и матричный вариант в \texttt{float64};
\item проверить инвариантность при $X_i\leftarrow X_i+c$ и
      $Y_i\leftarrow Y_i+c_y$;
\item сравнить разные размеры батча и перестановки соседей;
\item отдельно сравнить одну и ту же матрицу $W$ в плотном и CSR-форматах;
\item контролировать
\[
\eta_{jr}=
\frac{|\sum_iw_{ji}Q_{jri}|}
     {\sum_i|w_{ji}Q_{jri}|+\varepsilon};
\]
\item считать строки с малой массой, малым $n_{\mathrm{eff},j}$ или большим
      $\eta_{jr}$ численно ненадёжными.
\end{enumerate}
Для эталонных тестов следует использовать \texttt{float64}. На GPU допустимо хранение
в \texttt{float32}, но локальные массы, средние и редукции желательно аккумулировать
в повышенной точности. Формулы \eqref{eq:adp-I-stable},
\eqref{eq:adp-U-stable}, \eqref{eq:adp-sparse-I} и
\eqref{eq:adp-sparse-U} предпочтительнее некорректированного умножения $HX$.

\paragraph{Рекомендуемая реализация.}
Прямая версия формул \eqref{eq:adp-I-direct}--\eqref{eq:adp-U-direct}
используется как эталон для малых задач. Основное вычислительное ядро использует
факторизацию через $H$, произведение Адамара и матричные умножения. На CPU веса
хранятся в CSR, а на GPU --- в формате списка соседей. Центры и направления
обрабатываются батчами; данные предварительно центрируются и масштабируются, а
нулевой взвешенный момент проекций контролируется явно.

% Требуемые пакеты: amsmath, amssymb

\section{Матричное вычисление статистик Average Derivative Procedure}
\label{sec:adp-matrix-statistics}

Рассматриваются наблюдения
\[
X_i\in\mathbb R^d,\qquad Y_i\in\mathbb R,
\qquad i=1,\ldots,n,
\]
и набор локальных центров с индексами $j=1,\ldots,J$. Для каждого центра заданы веса
\[
w_{ji}=K\!\left(\left\|T(X_i-x_j)\right\|_2^2\right)\geq 0
\]
и $P=n_{\Phi}$ единичных направлений
$\phi_{jr}\in\mathbb R^d$, $r=1,\ldots,P$.
В ADP центрирование выполняется относительно локального взвешенного среднего
\begin{equation}
M_j=
\frac{\sum_{i=1}^n w_{ji}X_i}
     {\sum_{i=1}^n w_{ji}}.
\label{eq:adp-local-mean}
\end{equation}
Поэтому первый взвешенный момент равен нулю:
\begin{equation}
\sum_{i=1}^n w_{ji}(X_i-M_j)=0.
\label{eq:adp-zero-first-moment}
\end{equation}
Для каждого направления вычисляются статистики
\begin{align}
I_{jr}
&=
\sum_{i=1}^n
Y_i\,\phi_{jr}^{\mathsf T}(X_i-M_j)\,w_{ji},
\label{eq:adp-I-direct}
\\
U_{jr}
&=
\sum_{i=1}^n
(X_i-M_j)\,\phi_{jr}^{\mathsf T}(X_i-M_j)\,w_{ji}
\in\mathbb R^d.
\label{eq:adp-U-direct}
\end{align}
Строки $U_{jr}^{\mathsf T}$ образуют матрицу
$U_j\in\mathbb R^{P\times d}$, а значения $I_{jr}$ образуют
$I_j\in\mathbb R^P$. Локальная линейная модель приводит к приближению
\[
I_j\approx f'_jU_j\beta,
\qquad \|\beta\|_2=1.
\]
Свободный локальный коэффициент отсутствует именно благодаря
\eqref{eq:adp-zero-first-moment}.

\subsection{Плотная матричная факторизация}

Матрицу наблюдений будем хранить строками:
\[
X=
\begin{pmatrix}
X_1^{\mathsf T}\\[-1mm]
\vdots\\[-1mm]
X_n^{\mathsf T}
\end{pmatrix}
\in\mathbb R^{n\times d},
\qquad
y=(Y_1,\ldots,Y_n)^{\mathsf T}.
\]
Для батча из $B$ центров введём матрицу весов
$W\in\mathbb R^{B\times n}$ и вектор локальных масс
\begin{equation}
m=W\mathbf 1_n,
\qquad m_j=\sum_iw_{ji}.
\label{eq:adp-mass}
\end{equation}
Строки с $m_j\leq m_{\min}$ должны исключаться либо пересчитываться с большим bandwidth.
Нормированные веса задаются матрицей
\begin{equation}
A=\operatorname{diag}(m)^{-1}W,
\qquad A\mathbf 1_n=\mathbf 1_B.
\label{eq:adp-normalized-weights}
\end{equation}
Матрица локальных средних получается одним матричным умножением:
\begin{equation}
\mathcal M=AX\in\mathbb R^{B\times d},
\qquad \mathcal M_{j:}=M_j^{\mathsf T}.
\label{eq:adp-means-matmul}
\end{equation}

Направления для батча хранятся в тензоре
$\Phi\in\mathbb R^{B\times P\times d}$.
Сначала вычисляются проекции всех наблюдений:
\begin{equation}
R=\operatorname{bmm}(\Phi,X^{\mathsf T})
\in\mathbb R^{B\times P\times n},
\qquad
R_{jri}=\phi_{jr}^{\mathsf T}X_i.
\label{eq:adp-raw-projections}
\end{equation}
Проекции локальных средних имеют вид
\begin{equation}
c_{jr}=\phi_{jr}^{\mathsf T}M_j.
\label{eq:adp-projected-mean}
\end{equation}
После broadcasting получаем центрированные проекции
\begin{equation}
Q_{jri}=R_{jri}-c_{jr}
       =\phi_{jr}^{\mathsf T}(X_i-M_j).
\label{eq:adp-Q}
\end{equation}
Главный вспомогательный тензор строится произведением Адамара:
\begin{equation}
H=Q\odot A_{:\,,\mathrm{None},:},
\qquad
H_{jri}=a_{ji}Q_{jri}.
\label{eq:adp-H}
\end{equation}
Тогда обе статистики вычисляются через матричные произведения:
\begin{align}
I
&=
 m_{:\,,\mathrm{None}}\odot(Hy)
\in\mathbb R^{B\times P},
\label{eq:adp-I-matrix}
\\
U
&=
 m_{:\,,\mathrm{None},\mathrm{None}}
 \odot(HX)
\in\mathbb R^{B\times P\times d}.
\label{eq:adp-U-matrix}
\end{align}
Формула \eqref{eq:adp-U-matrix} следует из
\[
\sum_iw_{ji}Q_{jri}M_j
=M_j\phi_{jr}^{\mathsf T}\sum_iw_{ji}(X_i-M_j)=0.
\]
Следовательно, тензор разностей размера $B\times n\times d$ формировать не требуется.
Основной временный объект имеет размер $B\times P\times n$; при $P\ll d$
это существенно уменьшает память.

\paragraph{Устойчивая форма.}
Для защиты от ошибок округления рекомендуется предварительно глобально центрировать данные:
\[
X^{c}=X-\mathbf 1_nx_0^{\mathsf T},
\qquad
x_0=n^{-1}\sum_iX_i.
\]
Локальные средние вычисляются как $\mathcal M^{c}=AX^{c}$.
После построения $Q$ полезно повторно занулить его взвешенное среднее:
\begin{equation}
\delta_{jr}=\sum_i a_{ji}Q_{jri},
\qquad
Q_{jri}\leftarrow Q_{jri}-\delta_{jr}.
\label{eq:adp-Q-correction}
\end{equation}
Пусть $s=H\mathbf 1_n$ и $\bar y=Ay$. Численно защищённые формулы имеют вид
\begin{align}
I
&=
 m_{:\,,\mathrm{None}}
 \odot\left(Hy-s\odot\bar y\right),
\label{eq:adp-I-stable}
\\
U
&=
 m_{:\,,\mathrm{None},\mathrm{None}}
 \odot\left(
 HX^{c}-s_{:\,,\mathrm{None}}\odot
 \mathcal M^{c}_{:\,,\mathrm{None},:}
 \right).
\label{eq:adp-U-stable}
\end{align}
В точной арифметике $s=0$, поэтому
\eqref{eq:adp-I-stable}--\eqref{eq:adp-U-stable}
совпадают с \eqref{eq:adp-I-matrix}--\eqref{eq:adp-U-matrix}.
В плавающей арифметике дополнительные члены подавляют ошибку,
возникающую при умножении малого остатка $s$ на большой общий сдвиг $X$ или $y$.

\subsection{Батчевая схема}

Полную матрицу весов $J\times n$ и полный тензор проекций
$J\times P\times n$ хранить необязательно. Центры обрабатываются блоками
$j=j_0,\ldots,j_1-1$ размера $B$. Внутри каждого блока выполняются следующие шаги:
\begin{enumerate}
\item вычислить $W_b$, $m_b$, $A_b$ и $\mathcal M_b$;
\item вычислить $R_b=\operatorname{bmm}(\Phi_b,X^{\mathsf T})$;
\item получить $Q_b$, применить коррекцию \eqref{eq:adp-Q-correction};
\item выполнить $H_b=Q_b\odot A_b$;
\item вычислить $I_b$ и $U_b$ по
      \eqref{eq:adp-I-stable}--\eqref{eq:adp-U-stable};
\item записать результат и перейти к следующему блоку.
\end{enumerate}
Если число направлений велико, внутри блока центров вводится дополнительный блок
$r=r_0,\ldots,r_1-1$. Рабочая память плотной версии составляет
\[
O(BPn+BPd+Bn+Bd),
\]
а вычислительная сложность ---
\[
O(BPnd).
\]
Размер батча выбирается по памяти тензора $Q$ и результата $U$.

\subsection{Разреженная матрица весов}

Для ядра с компактным носителем большинство весов равно нулю.
Обозначим
\[
\mathcal N_j=\{i:w_{ji}\neq0\},
\qquad
E=\operatorname{nnz}(W),
\qquad
\bar k=E/J.
\]
На CPU матрицу $W$ целесообразно хранить в формате CSR.
Вектор масс и локальные средние вычисляются sparse--dense операциями:
\begin{equation}
m=W\mathbf 1_n,
\qquad
\mathcal M=\operatorname{diag}(m)^{-1}WX.
\label{eq:adp-sparse-mean}
\end{equation}
Для каждого ненулевого ребра $(j,i)$ и направления $r$ вычисляется только
\begin{equation}
q_{jri}=\phi_{jr}^{\mathsf T}(X_i-M_j),
\qquad i\in\mathcal N_j.
\label{eq:adp-sparse-q}
\end{equation}
Определим разреженную матрицу $H^{(r)}\in\mathbb R^{B\times n}$
с тем же паттерном ненулевых элементов, что у $W$:
\begin{equation}
H^{(r)}_{ji}=
\begin{cases}
\displaystyle \frac{w_{ji}}{m_j}q_{jri}, & i\in\mathcal N_j,\\[2mm]
0, & i\notin\mathcal N_j.
\end{cases}
\label{eq:adp-sparse-H}
\end{equation}
Тогда
\begin{align}
I_{:r}
&=
m\odot\left(H^{(r)}y-s^{(r)}\odot\bar y\right),
\label{eq:adp-sparse-I}
\\
U_{:r:}
&=
m_{:\,,\mathrm{None}}\odot
\left(
H^{(r)}X^{c}
-s^{(r)}_{:\,,\mathrm{None}}\odot\mathcal M^{c}
\right),
\label{eq:adp-sparse-U}
\\
s^{(r)}&=H^{(r)}\mathbf 1_n.
\end{align}
Операции $H^{(r)}y$ и $H^{(r)}X^{c}$ представляют соответственно
sparse matrix--vector и sparse matrix--dense matrix multiplication.
Не следует заранее хранить все $P$ матриц $H^{(r)}$: направления лучше
обрабатывать блоками, повторно используя один CSR-паттерн и меняя только массив значений.

Сложность разреженного варианта равна
\[
O(EPd)=O(J\bar kPd),
\]
а память для статистик и промежуточных значений ---
\[
O(EP+JPd+Jd).
\]
По сравнению с плотной схемой ожидаемое сокращение арифметики имеет порядок
$n/\bar k$.

\subsection{Формат списка соседей для GPU}

Если для каждого центра хранится не более $K$ соседей, разреженную структуру
удобно представить массивами
\[
N\in\mathbb N^{B\times K},
\qquad
W_{\mathrm{loc}}\in\mathbb R^{B\times K}.
\]
Здесь $N_{jk}$ --- индекс $k$-го соседа, а отсутствующие элементы маскируются
нулевым весом. После gather-операции
\[
X_{\mathcal N}=X[N]\in\mathbb R^{B\times K\times d},
\qquad
y_{\mathcal N}=y[N]\in\mathbb R^{B\times K}
\]
вычисления снова сводятся к batched matmul:
\begin{align}
M_j&=\sum_{k=1}^K a_{jk}X_{N_{jk}},
\\
Q_{jrk}&=\phi_{jr}^{\mathsf T}(X_{N_{jk}}-M_j),
\\
H_{jrk}&=a_{jk}Q_{jrk},
\\
I_{jr}&=m_j\sum_{k=1}^KH_{jrk}(Y_{N_{jk}}-\bar Y_j),
\\
U_{jr}&=m_j\sum_{k=1}^KH_{jrk}(X_{N_{jk}}-M_j).
\end{align}
Последние две редукции реализуются как
\[
(B,P,K)(B,K,1)\longrightarrow(B,P,1),
\qquad
(B,P,K)(B,K,d)\longrightarrow(B,P,d).
\]
Этот формат обычно предпочтительнее CSR на GPU, поскольку имеет регулярные размеры
и использует оптимизированные batched GEMM-ядра. Его недостаток --- padding при сильно
различающемся числе соседей.

\subsection{Численная устойчивость и контроль корректности}

Разреженность не меняет математическую задачу, если нулевые веса возникают из
компактного носителя ядра. Отсечение малых положительных весов или top-$K$ выбор
создают другую статистику. Для контроля следует вычислять долю отброшенной массы
\[
r_j=
\frac{\sum_{i:\,w_{ji}\ \mathrm{discarded}}w_{ji}}
     {\sum_iw_{ji}}
\]
и эффективное число соседей
\begin{equation}
n_{\mathrm{eff},j}
=
\frac{\left(\sum_iw_{ji}\right)^2}{\sum_iw_{ji}^2}
=
\frac{1}{\sum_i a_{ji}^2}.
\label{eq:adp-neff}
\end{equation}
Большое число ненулевых элементов не гарантирует устойчивость, если
$n_{\mathrm{eff},j}$ близко к единице.

Рекомендуемые проверки реализации:
\begin{enumerate}
\item сравнить прямые суммы и матричный вариант в \texttt{float64};
\item проверить инвариантность при $X_i\leftarrow X_i+c$ и
      $Y_i\leftarrow Y_i+c_y$;
\item сравнить разные размеры батча и перестановки соседей;
\item отдельно сравнить одну и ту же матрицу $W$ в плотном и CSR-форматах;
\item контролировать
\[
\eta_{jr}=
\frac{|\sum_iw_{ji}Q_{jri}|}
     {\sum_i|w_{ji}Q_{jri}|+\varepsilon};
\]
\item считать строки с малой массой, малым $n_{\mathrm{eff},j}$ или большим
      $\eta_{jr}$ численно ненадёжными.
\end{enumerate}
Для эталонных тестов следует использовать \texttt{float64}. На GPU допустимо хранение
в \texttt{float32}, но локальные массы, средние и редукции желательно аккумулировать
в повышенной точности. Формулы \eqref{eq:adp-I-stable},
\eqref{eq:adp-U-stable}, \eqref{eq:adp-sparse-I} и
\eqref{eq:adp-sparse-U} предпочтительнее некорректированного умножения $HX$.

\paragraph{Рекомендуемая реализация.}
Прямая версия формул \eqref{eq:adp-I-direct}--\eqref{eq:adp-U-direct}
используется как эталон для малых задач. Основное вычислительное ядро использует
факторизацию через $H$, произведение Адамара и матричные умножения. На CPU веса
хранятся в CSR, а на GPU --- в формате списка соседей. Центры и направления
обрабатываются батчами; данные предварительно центрируются и масштабируются, а
нулевой взвешенный момент проекций контролируется явно.
